\documentclass[../main.tex]{subfiles}
\begin{document}

\chapter*{How to Use This Guide}
\addcontentsline{toc}{chapter}{How to Use This Guide}
\markboth{How to Use This Guide}{}
\label{ch:howto}

This guide is built for the open-book exam, so it is engineered for fast retrieval rather than linear reading. It is split into two halves. The first half, the Overview at a Glance, is a compact map of the whole course. The second half, the numbered week chapters, is the detailed reference. The same four layers run through both halves, but they are deliberately separated so that scanning and deep reading do not get in each other's way.

\section*{The four layers}

The first layer is the \textbf{snapshot card} at the head of each week. It carries the weekly connection diagram, four to six exam-usable takeaways, and a Connects from and Connects to line that places the week inside the wider course flow.

The second layer is the \textbf{part summary}, one coloured box per lecture part. It states what the part establishes, why it matters, and where its full detail lives. It is the filter that answers whether the detail underneath is needed at all.

The third layer is the \textbf{bullet content} under each \emph{Content} heading. It reproduces the lecture logic faithfully, listing definitions, classifications, criteria, and the worked examples given in the slides. Terms defined for the first time are set in bold and registered in the index at the back.

The fourth layer is the \textbf{prose explanation} under each \emph{Explanation} heading. It turns the bullets into understanding by expanding the reasoning, the nuance, and the links between concepts.

\section*{The two halves}

The Overview at a Glance collects layers one and two for all six weeks back to back, with no detail in between. Read it to get the entire shape of the course in a few pages, or to locate which week and part a question belongs to. Each part summary there ends with a link to the exact section in the second half where its content lives.

The numbered week chapters hold layers three and four, the bullets and the prose, organised by part. Each chapter opens with a page reference back up to its overview entry, so navigation works in both directions. The cross-week threads, the report guide, and the appendices follow the week chapters.

\section*{The colour scheme}

The snapshot cards and headings follow the colours of the course connection diagrams. Green marks the conceptual track, the chain that runs from the language of research through conceptual design to the written report. Red marks the empirical track, the chain that runs from sampling and data through validation to data management. Blue-grey marks the integrative topics that draw both tracks together. The course map on the next page shows the two tracks in full.

\section*{Finding things on paper}

Every internal pointer uses \verb|\autoref|, which prints the entity type with the number, so a reference reads ``Section~2.3'' or ``Appendix~A'' rather than a bare number. This makes each pointer usable directly on a printed page. The fastest lookup tool is the index at the very back, where every bold term is listed with its page number. The glossary in Appendix~\ref{ch:glossary} gives one-line definitions with a pointer to the section where each term is treated in full.

\end{document}
